\documentclass[11pt]{article}

\usepackage[margin=0.85in]{geometry}
\usepackage[T1]{fontenc}
\usepackage[utf8]{inputenc}
\usepackage{amsmath,amssymb,bm}
\usepackage{booktabs,tabularx,array}
\usepackage{xcolor}
\usepackage{enumitem}
\usepackage{xurl}
\usepackage{hyperref}

\hypersetup{colorlinks=true,linkcolor=blue!50!black,urlcolor=blue!50!black}
\setlist{leftmargin=2em}
\setlength{\parskip}{0.25em}
\setlength{\parindent}{0pt}
\newcolumntype{Y}{>{\raggedright\arraybackslash}X}

\newcommand{\C}{\mathbb C}
\newcommand{\norm}[1]{\left\lVert #1\right\rVert}
\newcommand{\PSR}{\mathrm{PSR}}

\title{Number of Partial-FFT Views}
\author{}
\date{9 August 2026}

\begin{document}
\maketitle

\section{Introduction}

Underwater acoustic orthogonal frequency division multiplexing (OFDM)
receivers operate on channels that are time varying and frequency selective.
A receiver can demodulate one OFDM block with a single fast Fourier
transform (FFT) over the useful received interval.  When the channel changes
within that interval, the carriers mix, and the receiver sees intercarrier
interference (ICI).  One class of remedies keeps time-local components of
the interval.  The receiver transforms shorter parts of the interval
separately, so each transform sees the channel over a shorter time.
Partial-FFT demodulation is such a receiver \cite{yerramalli}.  It keeps
time-local components of the full-interval FFT, while combiner weights form
one pre-decoder soft symbol per carrier.

After demodulation, errors are corrected with forward error correction
(FEC).  The receivers in this note check parity constraints and reuse
decoder posterior information, so any linear block code with parity
constraints can serve.  Low-density parity-check (LDPC) codes are the common
choice, and every retained run uses LDPC as the demonstration.

We reconsider whether four partial-FFT views help at $N=1024$, why the
Partial-FFT receiver can lose to OFDM with FEC, and whether the result
depends on the code rate.  The pilot-trained combiner estimates its weights
from pilot carriers alone; we write pilot-only for the receiver that stops
there.  JUNA-Lite starts from that pilot-only candidate and refines it with
decoder posterior symbols as weighted anchors.  We separate three kinds of
statements: mathematical identities, measurements in the retained files,
and possible explanations that require another comparison.

All receivers at one printed configuration receive the same waveform.
Hence their uncoded and coded transmission rates are identical before receiver
errors are counted.  Receiver performance changes the effective data rate
through
\begin{equation}
  R_{\rm eff}
  =
  \frac{B_{\rm pay}}{T_{\rm frame}}\PSR,
  \label{eq:effective-rate}
\end{equation}
where $B_{\rm pay}$ is the payload bits per frame, $T_{\rm frame}$ is the
frame transmission duration, and $\PSR$ is the payload-exact frame success
rate.  This expression separates waveform rate from the fraction of frames
recovered by a receiver.  A receiver can therefore raise the delivered rate
only by recovering more frames.

The two result families answer different questions.  Table~\ref{tab:evidence}
records the boundary used throughout this note.

\begin{table}[ht]
\centering
\caption{Result families used in this note.}
\label{tab:evidence}
\begin{tabularx}{\textwidth}{@{}p{0.18\textwidth}YYY@{}}
\toprule
Result family & Channel and added noise & Frame use & Comparison supported \\
\midrule
First-four-second AWGN
& Measured Red replay with complex Gaussian noise that is independent over samples
& Four approximately non-overlapping frames per path and SNR, one seed
& Receiver counts at one fixed waveform with four views \\
Legacy added-noise sweep
& Measured Red replay with impulsive noise correlated across the three hydrophones
& Sixty frames per path and SNR, one seed; the same workloads recur across SNR after noise scaling
& Three matched directory pairs with one and four views at code rate 0.125 \\
Selected $N=1024$, CP$=128$ rate rows
& Legacy sweep with impulsive added noise
& Sixty frames per path and SNR, one seed
& Code rates 0.25 and 0.5 represented only by one-view runs in this selection \\
\bottomrule
\end{tabularx}
\end{table}

The first family is additive white Gaussian noise (AWGN), whereas the matched
view comparison is not.  We therefore do not use the legacy comparison to
identify the cause of the AWGN panel.  The historical rate-0.5 rows are not
part of the retained AWGN results.

The retained files support one description.  In every retained rate-0.125
comparison, OFDM+FEC results match exactly with one and four views, the
pilot-only receiver recovers fewer frames with four views, and JUNA-Lite
recovers more; the cause is not yet identified.  The remainder of this note
states what that sentence rests on, and which comparisons would identify the
cause.

\section{What the Number of Views Cannot Change}

We first state what changing the number of views cannot change.  Let $r[n]$
be the useful received interval and let $g_m[n]$ select view $m$.  We write
the carrier-$k$ view as
\begin{equation}
  Y_{m,k}
  =
  \sum_{n=0}^{N-1} g_m[n]r[n]e^{-j2\pi kn/N},
  \label{eq:view}
\end{equation}
where $N$ is the useful interval length, $n$ is the sample index, $k$ is the
carrier index, and $m$ is the view index.  Each view retains the contribution
from one non-overlapping part of the useful interval, so each view sees the
channel over a shorter time than the full transform does.

The implementation windows form a partition of the useful samples:
\begin{equation}
  \sum_{m=1}^{M}g_m[n]=1,
  \label{eq:partition}
\end{equation}
where $M$ is the number of partial-FFT views.  Linearity of the FFT then gives
\begin{equation}
  \sum_{m=1}^{M}Y_{m,k}=Y_k^{\rm full},
  \label{eq:sum-identity}
\end{equation}
where $Y_k^{\rm full}$ is the full-interval FFT output on carrier $k$.
Equation~\eqref{eq:sum-identity} holds for every number of non-overlapping
views used by this receiver.

The OFDM+FEC receiver applies residual pilot equalization after the sum in
\eqref{eq:sum-identity}.  It must therefore be independent of the configured
number of views, apart from numerical precision.  Across each of the three
legacy pairs, its successes, bit errors, bit error rate (BER), and effective
rate match exactly in all 192 path--SNR cells.  (The three pairs contain 576
such cells; none differs.)

This agreement is strong regression evidence, but the result file does not
prove the identity.  The identity follows from the partition and FFT
linearity; a numerical test protects its implementation.  The receivers can
therefore differ only after the transforms, in how the views are combined.
We describe that combiner next.

\section{The Pilot-Trained Combiner}

Partial-FFT reception is intended to retain information when the channel
varies within the full useful interval.  The pilot-only receiver therefore
does not use the sum with equal weights.  For carrier $k$ in band $b$, it
forms
\begin{equation}
  Z_k
  =
  \sum_{m=1}^{M}w_{m,b}Y_{m,k},
  \label{eq:combiner}
\end{equation}
where $w_{m,b}$ is the complex weight for view $m$ and carrier band $b$.
Unequal weights preserve a different combination of the time-local
components than the full FFT.

The weights solve a bandwise ridge problem:
\begin{equation}
  \widehat{\bm w}_b
  =
  \underset{\bm w\in\C^M}{\operatorname{arg\,min}}
  \left\{
  \sum_{p\in\mathcal P_b}
  \left|\bm w^{T}\bm Y_p-X_p\right|^2
  +\delta\norm{\bm w}_2^2
  \right\},
  \label{eq:ridge}
\end{equation}
where $\bm Y_p=[Y_{1,p},\ldots,Y_{M,p}]^T$, $\bm w$ is a trial direct
coefficient vector, $\widehat{\bm w}_b$ is the fitted vector, $X_p$ is the
known symbol, $\delta$ is the ridge term, and $\norm{\cdot}_2$ is the
Euclidean norm.  The set $\mathcal P_b$ contains the pilot targets in band
$b$; when the band contains fewer than $M$ targets, the implementation uses
all target pilots.  The transpose in \eqref{eq:ridge} matches the direct
coefficient sum in \eqref{eq:combiner}.  This fit trades pilot error against
the size of the combiner weights.

After combining, the receiver divides every pilot by its known symbol,
interpolates the resulting response over the active carriers, and divides
$Z_k$ by that response.  With one view and one global band, the fitted
combiner is one scalar $\alpha$, giving
\begin{equation}
  Z_k=\alpha Y_k^{\rm full},
  \qquad
  \widehat H_k=\alpha\widetilde H_k,
  \qquad
  \frac{Z_k}{\widehat H_k}
  =
  \frac{Y_k^{\rm full}}{\widetilde H_k},
  \label{eq:scalar-cancellation}
\end{equation}
where $\widetilde H_k$ is the response obtained by interpolating the unscaled
pilot responses.  The fitted scalar multiplies every carrier and every pilot
alike, so the division removes it.  The cancellation holds provided the
scalar does not reach the implementation's guard for a response near zero.

The deployed receiver uses 16 bands.  A different scalar can therefore be
fitted in each band, and pilot interpolation can cross a boundary between two
scalars.  With four views, the weighted result is generally not a scalar
multiple of the full FFT, so the cancellation in
\eqref{eq:scalar-cancellation} does not apply.

The current package test checks only that one raw view equals the full useful
interval FFT.  It does not check equality after combining one view with one
band spanning all active carriers and applying residual pilot equalization.
Thus the proposed PFFT-002 test is mathematically sound, but it is not yet
present in \texttt{Pkg.test()}.  We report the measured results next.

\section{Measured-Channel Results}

\subsection{The AWGN Panel}

We first examine the panel that motivated the question.  It uses $N=1024$,
CP$=64$, code rate 0.25, pilot spacings 5/5, check degree 14, and four views.
At 30~dB, pooling the twelve Red channel--hydrophone paths gives
Table~\ref{tab:awgn}.

\begin{table}[ht]
\centering
\caption{First-four-second AWGN frame counts at 30~dB.}
\label{tab:awgn}
\begin{tabular}{@{}lrr@{}}
\toprule
Receiver & Successful frames & Frame success rate \\
\midrule
OFDM+FEC & 42/48 & 0.875 \\
Pilot-only Partial-FFT+FEC & 36/48 & 0.750 \\
JUNA-Lite & 46/48 & 0.958 \\
JUNA $(C,z)$ experiment receiver & 46/48 & 0.958 \\
JUNA $(C,W,z)$ experiment receiver & 46/48 & 0.958 \\
\bottomrule
\end{tabular}
\end{table}

The count of four frames per path is too small to order the receivers.
The OFDM and JUNA-Lite totals differ by four frames, and the saved files do not
retain their paired frame outcomes.  We therefore report the counts without
claiming an improvement for receivers using the same waveform.

\subsection{The Matched View Pairs}

The matched legacy comparison is larger but belongs to the impulsive added-noise
family.  Table~\ref{tab:views} pools SNR from 20 to 30~dB, giving 4320
frame trials represented by aggregate counts per receiver and column.

\begin{table}[ht]
\centering
\caption{Legacy frame success rates with impulsive added noise and four or one partial-FFT views.}
\label{tab:views}
\begin{tabular}{@{}llrr@{}}
\toprule
Geometry & Receiver & Four views & One view \\
\midrule
CP16, pilots 3/5
& OFDM+FEC & 0.827778 & 0.827778 \\
& Pilot-only Partial-FFT+FEC & 0.566435 & 0.816204 \\
& JUNA-Lite & 0.930787 & 0.815509 \\
\addlinespace
CP16, pilots 3/10
& OFDM+FEC & 0.811343 & 0.811343 \\
& Pilot-only Partial-FFT+FEC & 0.546296 & 0.802083 \\
& JUNA-Lite & 0.931481 & 0.803241 \\
\addlinespace
CP128, pilots 5/10
& OFDM+FEC & 0.698843 & 0.698843 \\
& Pilot-only Partial-FFT+FEC & 0.371991 & 0.674769 \\
& JUNA-Lite & 0.896991 & 0.677315 \\
\bottomrule
\end{tabular}
\end{table}

Four views reduce the pilot-only success rate in all three pairs, from 0.82
to 0.57, from 0.80 to 0.55, and from 0.67 to 0.37.  The same four views
raise the JUNA-Lite rate in all three pairs, from 0.82 to 0.93, from 0.80 to
0.93, and from 0.68 to 0.90.  OFDM+FEC matches exactly in every cell.  The
same decomposition therefore harms the receiver that trains its weights on
pilots alone and helps the receiver that refines with decoder posterior
symbols.

The $(C,z)$ and $(C,W,z)$ experiment receivers differ from JUNA-Lite by no
more than 0.0031 in these comparisons with four views.  The gradient-only
receivers therefore do not materially change the recorded pooled rates; the
large increase is already present in JUNA-Lite.

The table does not constitute a four-view comparison with equal weights.
Changing from one to four views changes the dimension of the ridge fit, the
view signals, and the stopping behavior of JUNA-Lite.  A receiver with four
views and one common fixed weight is required to isolate the decomposition
itself.

\subsection{The Code-Rate Rows}

The available rows do not separate code rate from the number of views.
Table~\ref{tab:rate} shows the settings behind the earlier interpretation that
the result depends on code rate.

\begin{table}[ht]
\centering
\caption{Historical high-SNR comparison at $N=1024$ and CP$=128$.}
\label{tab:rate}
\begin{tabular}{@{}rrrrl@{}}
\toprule
Code rate & Views & OFDM+FEC & Best JUNA receiver & Reading \\
\midrule
0.125 & 4 & 0.698843 & 0.897917 & Four-view result \\
0.25  & 1 & 0.735880 & 0.725926 & One-view result \\
0.5   & 1 & 0.176620 & 0.162269 & One-view legacy result \\
\bottomrule
\end{tabular}
\end{table}

The table supports a gain with four views in the tested rate-0.125 legacy
configuration.  It does not show that the gain is confined to that code rate,
because the other code rates were measured with one view.  Failure to measure
a difference is also not evidence that two receivers are equivalent.  Before
interpreting these differences, we state what the retained files can and
cannot support.

\section{Limitations}

Every receiver within one run receives a copy of the same received waveform,
so the receiver outcomes are paired.  The calculation that treats the two receiver
proportions as independent gives $z=1.477$ and $p=0.140$ for the AWGN counts.
However, that calculation does not match
the experiment.  From the two margins alone, the possible exact two-sided
McNemar $p$-values range from 0.125 to 0.289.

The legacy calculations that give $p$-values below $10^{-50}$ also treat
4320 outcomes per receiver as independent.  However, each SNR uses the same sixty
frame seeds, the receivers share each received waveform, and the three hydrophones
share correlated added noise.  With four capture groups, the minimum
attainable two-sided exact sign-test $p$-value is 0.125.

We therefore use the legacy differences as descriptions of these retained
files.  Repeated seeds support statements about these captures, while more
independent captures are required for statements beyond the retained Red
files.  Both cases require frame traces for every receiver and a paired comparison
that keeps the capture as a group.  Overlapping marginal intervals do not
replace that paired comparison.  With those limits stated, we turn to the
possible explanations.

\section{Possible Explanations}

The retained results identify several effects but do not select one cause.
We therefore state each explanation with the measurement that would oppose
it.

\subsection{Pilots per Band}

At $N=1024$, pilot spacing 5 gives 205 pilots, or only 12--13 pilots in each
of 16 bands.  The pilot-only receiver estimates four complex weights per band
from these targets.  At $N=2048$ and $N=4096$, the corresponding counts are
25--26 and 51--52 pilots per band.

The four columns of the ridge problem can be correlated because every view
comes from the same received block.  A poorly conditioned correlation matrix
can make the fitted weights sensitive to pilot noise, even when the ridge term
keeps the solve finite.  This explanation predicts large condition numbers and
large weight norms where pilot-only Partial-FFT reception loses.

For uncorrelated view noise, the combined noise variance is
\begin{equation}
  \sigma_Z^2
  =
  \sum_{m=1}^{M}|w_{m,b}|^2\sigma_m^2,
  \label{eq:noise-gain}
\end{equation}
where $\sigma_m^2$ is the noise variance in view $m$ and $\sigma_Z^2$ is the
variance after combining.  Large fitted weights can therefore increase the
noise passed to the decoder.

The implementation did not retain the weights, correlation matrix condition
numbers, or the noise gain.  The explanation based on the number of pilots is
therefore consistent with the measured $N$ crossover --- in the AWGN family
the pilot-only receiver loses at $N=1024$ and wins at $N=2048$ and $N=4096$
--- but the current files do not establish it.

\subsection{Band Boundaries}

One weight vector is shared by every carrier in a band.  Residual pilot
equalization then interpolates pilot responses across active carriers,
including carriers near a band boundary.  Different weight vectors on the two
sides need not pass through that interpolation as one common factor.

This explanation predicts a smaller gap with one band spanning all carriers, with weights
that vary continuously across frequency, or with the same fixed weights in
all bands.  It also explains why the identity for one view and one band spanning
all carriers says
nothing about the deployed 16-band path.

\subsection{OFDM Under Residual Doppler}

Shorter FFT intervals reduce mixing among carriers, and weighted combining of
their outputs can reduce intercarrier interference
\cite{yerramalli}.  In the notation of this note, equal weights reconstruct
the full FFT, while unequal weights retain a different combination of the
view coefficients.

At a sampling rate of 9.6~kHz, the useful interval and one of four views have
the durations in Table~\ref{tab:duration}.

\begin{table}[ht]
\centering
\caption{Useful-interval and single-view durations for four views at 9.6~kHz.}
\label{tab:duration}
\begin{tabular}{@{}rrr@{}}
\toprule
$N$ & Useful interval & One of four views \\
\midrule
1024 & 106.7~ms & 26.7~ms \\
2048 & 213.3~ms & 53.3~ms \\
4096 & 426.7~ms & 106.7~ms \\
\bottomrule
\end{tabular}
\end{table}

On a time-varying channel, longer OFDM blocks can expose more within-block
variation to the views.  However,
larger $N$ also supplies more pilots per band, so the measured crossover
cannot be assigned to block duration alone.  A comparison with the channel
held fixed and another with the same number of pilots are both required.

\subsection{Candidate Selection and Stopping}

The frame-wide JUNA-Lite path starts from the pilot-only Partial-FFT candidate.
It returns immediately when that candidate satisfies the low-density
parity-check syndrome.  In this path, the validity flag means a zero
syndrome and does not mean that a cyclic redundancy check (CRC) passed.

Four-view pilot-only reception performs poorly in the three legacy pairs, so
it can cause JUNA-Lite to enter posterior-anchor refinement more often.  The
pilot-only reception with one view performs much better and can stop refinement
earlier.  The measured JUNA-Lite difference may therefore combine useful
time-local information with a change in how often refinement runs.

This explanation predicts that the gap changes when both view settings run
the same fixed number of refinements.  It also predicts parity-valid but
CRC-invalid early returns if a wrong codeword satisfies the parity checks.
The saved Lite files do not retain these events.

\subsection{LDPC Posterior Symbols}

A lower code rate supplies more redundant bits and can give the decoder more
reliable posterior information for the JUNA-Lite anchor set.  Decoder feedback
is useful when pilots alone do not specify stable combiner weights and when the
posterior information is reliable.

This explanation predicts a larger Lite gain with four views at lower code rate
after the number of views, noise realization, geometry, and stopping rule are
held fixed.  The existing rate table does not make that comparison, so the
explanation based on code rate remains untested.

\subsection{Cyclic Prefix}

The cyclic prefix (CP) reduces intersymbol interference when it covers the
channel delay spread.  For each path in the audit, squared tap magnitudes were
averaged over replay snapshots to form a mean power-delay profile.  The
shortest contiguous interval containing 90 percent of its energy ranges from
about 3.1 to 8.1~ms.  CP64 is 6.67~ms and covers eleven of the twelve mean
windows under that rule, while CP128 is 13.33~ms and covers all twelve.

The 90-percent rule is a design choice rather than a physical boundary.  A
shorter CP leaves residual energy outside the guard, while a longer CP spends
more transmission time on the guard.  CP128 is therefore a conservative guard
under this rule, not the only usable guard.  The audit calculation does not yet
have a tracked script and is not a reproducible result until that script is
retained.

A coherence calculation using the tap estimates alone also omits the separate
samplewise phase trajectory used by the replay channel.  The quoted coherence time therefore
cannot be used as a hard limit that discards every longer-$N$ result.  Four-view
reception is intended to retain information when the channel varies within the
full useful interval.

\section{Complexity and Evaluation}

The simple OFDM+FEC receiver does not run the Partial-FFT combiner refinement.
Both receivers can stop after one decoder run when the initial JUNA-Lite
candidate is valid.  JUNA-Lite adds combiner fits and decoder runs only when
refinement proceeds.  At one waveform, it cannot increase the transmitted raw
rate; it can increase the delivered rate only by recovering more frames.

The present Red evidence does not support the earlier statement that JUNA
should not be used.  It supports a narrower statement: JUNA-Lite with four
views has a large descriptive gain in three rate-0.125 legacy comparisons
with impulsive added noise, while the AWGN comparison with four frames per
path remains inconclusive.

One proposed policy would decode with Standard first, let a failed CRC trigger
JUNA-Lite, use CRC for Lite stopping and final selection, and return an erasure
if Lite also fails.  That policy was not measured.  The AWGN $(C,z)$ receivers
set \texttt{cz\_crc\_gate=false} and \texttt{cz\_gradient\_only=true}, while
frame-wide JUNA-Lite starts from pilot-only Partial-FFT reception and stops on
LDPC validity rather than CRC validity.  The measured receivers do not
demonstrate a CRC-gated policy.

Until the controlled comparisons run, the retained files support one working
rule for these captures at code rate 0.125.  One should change the number of
views only together with the receiver that uses them: with pilots alone, one
view recovered more frames; with posterior-anchor refinement, four views
recovered more.  The cost of the four-view choice is the refinement's
combiner fits and decoder runs.  The choice is undone by returning the view
count to one, which leaves OFDM+FEC unchanged by the identity in
\eqref{eq:sum-identity}.  The comparisons that would settle the remaining
questions follow.

\section{Experimental Method}

We can distinguish the explanations with the controlled comparisons in
Table~\ref{tab:tests}.  Each comparison keeps the received frames, pilot
symbols, code, channel, and added noise fixed unless the named quantity is the
one being changed.

\begin{table}[ht]
\centering
\caption{Controlled comparisons for the remaining explanations.}
\label{tab:tests}
\begin{tabularx}{\textwidth}{@{}p{0.23\textwidth}YY@{}}
\toprule
Question & Controlled comparison & Evidence to retain \\
\midrule
One-view cancellation
& One view, one global band, and the same input through Standard and Partial-FFT equalization
& Equalized arrays before decoding, with a numerical tolerance \\
Four-view reconstruction
& Four views with every weight fixed to the same global value, against Standard
& Equalized arrays, BER, and frame success \\
Band-boundary effect
& One band against 16 bands with the same four views and targets
& Per-band weights and carrier error near each boundary \\
Weight-estimation loss
& Pilot-only reception with four views and logged ridge matrices
& Condition numbers, weight norms, pilot residuals, and noise gain \\
Lite stopping effect
& One and four views with the same fixed number of Lite refinements
& Initial syndrome, CRC, selected iteration, and every candidate outcome \\
Code-rate effect
& Matched AWGN runs with one and four views at rates 0.125 and 0.25
& Paired frame outcomes for every receiver; rate 0.5 only if that experiment is reopened \\
Channel-variation effect
& Frozen taps and phase against the measured replay, with pilot support held fixed
& Effective-channel variation, weights, BER, and frame success \\
Deployment policy
& Standard first, followed by CRC-gated Lite only after failure
& Effective data rate, processing latency, decoder runs, and undetected errors \\
\bottomrule
\end{tabularx}
\end{table}

The frame traces must retain Standard, pilot-only Partial-FFT, and JUNA-Lite,
not only the gradient-only receivers.  Repeated seeds and additional
independent captures are needed before probabilities are used to generalize
beyond the retained Red files.

We record the branches, exact experiment directories, result manifests, and
receiver files so that each statement can be checked.

\begingroup
\small
\begin{itemize}[itemsep=0.15em,topsep=0.2em]
  \item AWGN panel: branch \texttt{agent/awgn-results}, decision AWGN-018,
  and manifest
  \path{JunaCore/experiments/2026-08-08-red-awgn-first4s-frames4-snr-sweep-n1024-cp64-rate025-p5-5-dc14-kfill-pfft4/results/results_manifest.json}.
  Its configuration is $N=1024$, CP$=64$, rate 0.25, pilots 5/5, check
  degree 14, and four views.  The manifest lists the twelve configuration
  CSV files used for Table~\ref{tab:awgn}.
  \item Legacy matched pairs: branch \texttt{agent/all-results}.  Under
  \path{JunaCore/experiments}, the exact directory pairs are
  \begin{itemize}
    \item \path{2026-08-06-red-snr-sweep-n1024-cp16-rate0125-p3-5-dc10-kfill}
    and \path{2026-08-06-red-snr-sweep-n1024-cp16-rate0125-p3-5-dc10-kfill-pfft1};
    \item \path{2026-08-06-red-snr-sweep-n1024-cp16-rate0125-p3-10-dc10-kfill}
    and \path{2026-08-06-red-snr-sweep-n1024-cp16-rate0125-p3-10-dc10-kfill-pfft1};
    \item \path{2026-08-06-red-snr-sweep-n1024-cp128-rate0125-p5-10-dc10-kfill}
    and \path{2026-08-06-red-snr-sweep-n1024-cp128-rate0125-p5-10-dc10-kfill-pfft1}.
  \end{itemize}
  Each directory contains \path{results/results_manifest.json}; each manifest
  lists the twelve configuration CSV files used for Table~\ref{tab:views}.
  \item AWGN receiver provenance: the result manifest pins these values:
  \begin{itemize}
    \item source commit:\\
    \path{14976b6c057549dcc00b5cd5020a97d6cfc7ffdb}
    \item receiver-source tree:\\
    \path{2db4aadfbbe9bde6cf34cb20e98a638d0f9aad0f}
  \end{itemize}
  \item Restored receiver source inspected for this note: main commit
  \texttt{7fbdec9}, especially
  \path{JunaCore/src/juna/common.jl} and
  \path{JunaCore/src/juna/frame_wide_ldpc.jl}.
  \item Existing regression test:
  \path{JunaCore/test/partial_fft_receiver.jl}.
\end{itemize}
\endgroup

The older legacy manifests do not record a source commit or tree, and the
default files with four views do not print the view count in their CSV schema.
Their directory pairs differ in the run wrapper by the explicit one-view
setting, and OFDM results match exactly, but the missing provenance remains a
reproducibility limitation.

\section{Conclusion}

In every retained rate-0.125 comparison, OFDM+FEC results match exactly with
one and four views, the pilot-only receiver recovers fewer frames with four
views, and JUNA-Lite recovers more; the cause is not yet identified.

We considered why four partial-FFT views can reduce pilot-only receiver
performance at $N=1024$ while improving JUNA-Lite.  Summation with equal weights
reconstructs the full FFT for every view count, but pilot-only Partial-FFT
reception estimates unequal bandwise weights from limited pilots.  JUNA-Lite
then adds decoder posterior symbols as weighted anchors, unless a valid initial
candidate stops the refinement.

The retained legacy results show a large JUNA-Lite increase with four views at code
rate 0.125 and a large pilot-only decrease under impulsive added noise.  They
do not establish dependence on code rate, identify one physical cause, or
validate the equality test with one view and one band spanning all carriers.
The AWGN panel with four frames per path is also too small to order the receivers.

Under the stated 90-percent delay-window rule, CP128 is a conservative guard,
while $N=1024$ remains a tested choice for the simple OFDM receiver.  The tap
estimates alone do not establish that the replay channel stays fixed over the
$N=1024$ useful interval.
The retained files therefore do not justify excluding JUNA-Lite.  They show a
descriptive four-view increase in the legacy rate-0.125 cases.  Its AWGN
improvement, code-rate dependence, cause, and processing trade-off remain
undetermined until the controlled comparisons in Table~\ref{tab:tests} use
matched frames and settings.

\begin{thebibliography}{9}

\bibitem{yerramalli}
S. Yerramalli, M. Stojanovic, and U. Mitra,
``Partial FFT Demodulation: A Detection Method for Highly Doppler Distorted
OFDM Systems,''
\emph{IEEE Transactions on Signal Processing}, vol.~60, no.~11,
pp.~5906--5918, November 2012.

\end{thebibliography}

\end{document}
